% ch7_a 微扰理论与引力辐射 (书页 274-286 / p289-p301)
% 第7章 Perturbation Theory and Gravitational Radiation
\chapter{微扰理论与引力辐射}

\section{线性化引力与规范变换}

% ---- 书页 274 / p289 ----
我们最初推导 Einstein 方程的时候，是通过考虑 Newton 极限来确认自己走在正确的道路上。当时我们对这句话的理解是：不仅引力场是弱的，而且它是静态的（不含时间导数），检验粒子也在缓慢地运动。本章要描述的弱场极限限制要少一些：只假定场仍然是弱的，但它可以随时间变化，并且对检验粒子的运动不加任何限制。这将使我们能够讨论一些在 Newton 理论中缺席或含混的现象，例如引力辐射（场随时间变化）以及光线偏折（涉及快速运动的粒子）。

引力场的微弱，再一次表现为我们有能力把度规分解成平直的 Minkowski 度规加上一个小微扰，
\begin{equation*}
g_{\mu\nu} = \eta_{\mu\nu} + h_{\mu\nu}, \qquad |h_{\mu\nu}| \ll 1.
\tag{7.1}
\end{equation*}
我们将把自己限制在 $\eta_{\mu\nu}$ 取规范形式的坐标系中，即 $\eta_{\mu\nu} = \mathrm{diag}(-1,+1,+1,+1)$。$h_{\mu\nu}$ 是小量这一假设，使我们能够忽略任何高于该量一阶的东西，由此立即得到
\begin{equation*}
g^{\mu\nu} = \eta^{\mu\nu} - h^{\mu\nu},
\tag{7.2}
\end{equation*}
其中 $h^{\mu\nu} = \eta^{\mu\rho}\eta^{\nu\sigma}h_{\rho\sigma}$。与之前一样，我们可以用 $\eta^{\mu\nu}$ 和 $\eta_{\mu\nu}$ 升降指标，因为相应的修正将是微扰的高阶小量。事实上，我们可以把广义相对论的线性化版本（其中 $h_{\mu\nu}$ 的高于一阶的效应都被略去）看作是在平直背景时空上传播的一个对称张量场 $h_{\mu\nu}$ 的理论。这个理论在狭义相对论的意义下是洛伦兹不变的；在洛伦兹变换 $x^{\mu'} = \Lambda^{\mu'}{}_{\mu}x^{\mu}$ 之下，平直度规 $\eta_{\mu\nu}$ 不变，而微扰按下式变换
\begin{equation*}
h_{\mu'\nu'} = \Lambda^{\mu'}{}_{\mu}\Lambda^{\nu'}{}_{\nu}h_{\mu\nu}.
\tag{7.3}
\end{equation*}

% ---- 书页 275 / p290 ----
注意，我们本来也可以考虑 Minkowski 空间之外的某个其他背景时空上的小微扰。在那种情形下，度规将被写成 $g_{\mu\nu} = g_{\mu\nu}^{(0)} + h_{\mu\nu}$，而我们导出的将是关于一个在度规为 $g_{\mu\nu}^{(0)}$ 的弯曲空间上传播的对称张量的理论。例如在宇宙学中，这样的处理方式就是必需的。

我们想找出微扰 $h_{\mu\nu}$ 所服从的运动方程，它来自把 Einstein 方程考察到一阶。我们从 Christoffel 符号入手，它由下式给出
\begin{align*}
\Gamma^{\rho}_{\mu\nu} &= \frac{1}{2}g^{\rho\lambda}(\partial_{\mu}g_{\nu\lambda} + \partial_{\nu}g_{\lambda\mu} - \partial_{\lambda}g_{\mu\nu}) \\
&= \frac{1}{2}\eta^{\rho\lambda}(\partial_{\mu}h_{\nu\lambda} + \partial_{\nu}h_{\lambda\mu} - \partial_{\lambda}h_{\mu\nu}).
\tag{7.4}
\end{align*}
由于联络系数是一阶量，它们对 Riemann 张量的贡献将只来自 $\Gamma$ 的导数，而不含 $\Gamma^{2}$ 项。为方便起见降低一个指标，我们得到
\begin{align*}
R_{\mu\nu\rho\sigma} &= \eta_{\mu\lambda}\partial_{\rho}\Gamma^{\lambda}_{\nu\sigma} - \eta_{\mu\lambda}\partial_{\sigma}\Gamma^{\lambda}_{\nu\rho} \\
&= \frac{1}{2}(\partial_{\rho}\partial_{\nu}h_{\mu\sigma} + \partial_{\sigma}\partial_{\mu}h_{\nu\rho} - \partial_{\sigma}\partial_{\nu}h_{\mu\rho} - \partial_{\rho}\partial_{\mu}h_{\nu\sigma}).
\tag{7.5}
\end{align*}
对 $\mu$ 和 $\rho$ 缩并便得到 Ricci 张量，
\begin{equation*}
R_{\mu\nu} = \frac{1}{2}(\partial_{\sigma}\partial_{\nu}h^{\sigma}{}_{\mu} + \partial_{\sigma}\partial_{\mu}h^{\sigma}{}_{\nu} - \partial_{\mu}\partial_{\nu}h - \Box h_{\mu\nu}),
\tag{7.6}
\end{equation*}
它显然关于 $\mu$ 和 $\nu$ 对称。在这个表达式中，我们把微扰的迹定义为 $h = \eta^{\mu\nu}h_{\mu\nu} = h^{\mu}{}_{\mu}$，而 d'Alembert 算符就是平直空间中的那一个：$\Box = -\partial_{t}^{2} + \partial_{x}^{2} + \partial_{y}^{2} + \partial_{z}^{2}$。再一次缩并以得到 Ricci 标量，就有
\begin{equation*}
R = \partial_{\mu}\partial_{\nu}h^{\mu\nu} - \Box h.
\tag{7.7}
\end{equation*}
把这些都放在一起，我们就得到了 Einstein 张量：
\begin{align*}
G_{\mu\nu} &= R_{\mu\nu} - \frac{1}{2}\eta_{\mu\nu}R \\
&= \frac{1}{2}(\partial_{\sigma}\partial_{\nu}h^{\sigma}{}_{\mu} + \partial_{\sigma}\partial_{\mu}h^{\sigma}{}_{\nu} - \partial_{\mu}\partial_{\nu}h - \Box h_{\mu\nu} - \eta_{\mu\nu}\partial_{\rho}\partial_{\lambda}h^{\rho\lambda} + \eta_{\mu\nu}\Box h).
\tag{7.8}
\end{align*}
与我们把线性化理论理解为描述平直背景上的对称张量的理论相一致，线性化 Einstein 张量 (7.8) 可以通过对下面这个拉格朗日量关于 $h_{\mu\nu}$ 作变分而导出：
\begin{align*}
\mathcal{L} = \frac{1}{2}\Big[(\partial_{\mu}h^{\mu\nu})(\partial_{\nu}h) &- (\partial_{\mu}h^{\rho\sigma})(\partial_{\rho}h^{\mu}{}_{\sigma}) + \frac{1}{2}\eta^{\mu\nu}(\partial_{\mu}h^{\rho\sigma})(\partial_{\nu}h_{\rho\sigma}) \\
&- \frac{1}{2}\eta^{\mu\nu}(\partial_{\mu}h)(\partial_{\nu}h)\Big].
\tag{7.9}
\end{align*}
这个拉格朗日量是否恰当，请你在习题中加以验证。

% ---- 书页 276 / p291 ----
线性化场方程当然就是 $G_{\mu\nu} = 8\pi GT_{\mu\nu}$，其中 $G_{\mu\nu}$ 由式 (7.8) 给出，而 $T_{\mu\nu}$ 是能量-动量张量，按 $h_{\mu\nu}$ 的零阶计算。我们不把更高阶的修正包括进能量-动量张量，因为要使弱场极限适用，能量和动量的总量本身就必须很小。换句话说，$T_{\mu\nu}$ 中最低的非零阶自动与微扰具有相同的数量级。注意，最低阶下的守恒定律就是简单的 $\partial_{\mu}T^{\mu\nu} = 0$。我们常常关心的是真空方程，它一如既往就只是 $R_{\mu\nu} = 0$，其中 $R_{\mu\nu}$ 由式 (7.6) 给出。

有了线性化场方程，我们差不多就准备好着手求解它了。然而，我们首先应该处理规范不变性这个棘手的问题。这个问题之所以出现，是因为要求 $g_{\mu\nu} = \eta_{\mu\nu} + h_{\mu\nu}$ 并没有完全确定时空上的坐标系；可能存在别的坐标系，度规在其中仍然可以写成 Minkowski 度规加上一个小微扰，只不过微扰有所不同。因此，把度规分解成平直背景加微扰的做法并不是唯一的。为了考察这个问题，我们将利用附录 A 和附录 B 中讨论的关于微分同胚的思想；尚未读过那两节的读者可以跳到式 (7.14) 及其后的两段，那里包含了最关键的思想。

让我们从一种“高观点”（highbrow）的角度来思考规范不变性。线性化理论可以看作是一个支配平直背景上张量场行为的理论——这个观念可以借助一个\textbf{背景时空}（background spacetime）$M_{b}$、一个\textbf{物理时空}（physical spacetime）$M_{p}$ 以及一个微分同胚 $\phi : M_{b} \to M_{p}$ 来形式化。作为流形，$M_{b}$ 和 $M_{p}$ 是同一个（因为它们微分同胚），但我们设想它们拥有一些不同的张量场：在 $M_{b}$ 上我们定义了平直的 Minkowski 度规 $\eta_{\mu\nu}$，而在 $M_{p}$ 上有某个遵守 Einstein 方程的度规 $g_{\alpha\beta}$。（我们设想 $M_{b}$ 配有坐标 $x^{\mu}$，$M_{p}$ 配有坐标 $y^{\alpha}$，尽管这些坐标并不会扮演什么显要角色。）微分同胚 $\phi$ 使我们能够在背景时空与物理时空之间来回搬移张量，如图 7.1 所示。既然我们想把线性化理论构造成一个发生在平直背景时空上的理论，我们感兴趣的便是物理度规的拉回 $(\phi^{*}g)_{\mu\nu}$。我们可以把微扰定义为拉回后的物理度规与平直度规之差：

%FIG{7.1}: {一个微分同胚把背景时空 $M_{b}$（具有平直度规 $\eta_{\mu\nu}$）与物理时空 $M_{p}$ 联系起来。}

% ---- 书页 277 / p292 ----
%FIG{7.2}: {由背景时空 $M_{b}$ 上的矢量场 $\xi^{\mu}$ 生成的单参数微分同胚族 $\psi_{\epsilon}$。}
\begin{equation*}
h_{\mu\nu} = (\phi^{*}g)_{\mu\nu} - \eta_{\mu\nu}.
\tag{7.10}
\end{equation*}
根据这个定义，$h_{\mu\nu}$ 的分量没有理由必须是小的；不过，如果 $M_{p}$ 上的引力场是弱的，那么对某些微分同胚 $\phi$ 我们就会有 $|h_{\mu\nu}| \ll 1$。于是我们只把注意力限制在那些使这一点成立的微分同胚上。这样一来，$g_{\alpha\beta}$ 在物理时空上遵守 Einstein 方程这一事实，就意味着 $h_{\mu\nu}$ 在背景时空上将遵守线性化方程（因为 $\phi$ 作为微分同胚，同样可以用来拉回 Einstein 方程本身）。

在这种语言里，规范不变性问题不过是说，在 $M_{b}$ 与 $M_{p}$ 之间存在着大量“许可的”微分同胚（“许可”的意思是微扰是小的）。考虑背景时空上的一个矢量场 $\xi^{\mu}(x)$。这个矢量场生成一个单参数微分同胚族 $\psi_{\epsilon} : M_{b} \to M_{b}$，如图 7.2 所示。当 $\epsilon$ 足够小时，如果 $\phi$ 是一个使由式 (7.10) 定义的微扰很小的微分同胚，那么 $(\phi \circ \psi_{\epsilon})$ 也是，只不过微扰将取不同的值。具体地，我们可以定义一个以 $\epsilon$ 为参数的微扰族：
\begin{align*}
h^{(\epsilon)}_{\mu\nu} &= [(\phi \circ \psi_{\epsilon})^{*}g]_{\mu\nu} - \eta_{\mu\nu} \\
&= [\psi_{\epsilon}^{*}(\phi^{*}g)]_{\mu\nu} - \eta_{\mu\nu}.
\tag{7.11}
\end{align*}
第二个等号基于这样的事实：复合映射的拉回等于各拉回按相反次序的复合，而这又源于拉回本身就是沿与原映射相反的方向搬动事物。代入关系式 (7.10)，我们发现
\begin{align*}
h^{(\epsilon)}_{\mu\nu} &= \psi_{\epsilon}^{*}(h + \eta)_{\mu\nu} - \eta_{\mu\nu} \\
&= \psi_{\epsilon}^{*}(h_{\mu\nu}) + \psi_{\epsilon}^{*}(\eta_{\mu\nu}) - \eta_{\mu\nu},
\tag{7.12}
\end{align*}
这是因为两个张量之和的拉回等于各自拉回之和。现在我们动用 $\epsilon$ 是小量这一假设；此时 $\psi_{\epsilon}^{*}(h_{\mu\nu})$ 在最低阶下等于 $h_{\mu\nu}$，而另外两项则给我们一个李导数：

% ---- 书页 278 / p293 ----
\begin{align*}
h^{(\epsilon)}_{\mu\nu} &= \psi_{\epsilon}^{*}(h_{\mu\nu}) + \epsilon\left[\frac{\psi_{\epsilon}^{*}(\eta_{\mu\nu}) - \eta_{\mu\nu}}{\epsilon}\right] \\
&= h_{\mu\nu} + \epsilon\mathcal{L}_{\xi}\eta_{\mu\nu}.
\tag{7.13}
\end{align*}
在附录 B 中我们证明过，度规沿矢量场 $\xi_{\mu}$ 的李导数为 $\mathcal{L}_{\xi}g_{\mu\nu} = 2\nabla_{(\mu}\xi_{\nu)}$。在当前的情形下，背景度规是平直的，协变导数变成偏导数；因此我们有
\begin{equation*}
\boxed{\,h^{(\epsilon)}_{\mu\nu} = h_{\mu\nu} + 2\epsilon\partial_{(\mu}\xi_{\nu)}.\,}
\tag{7.14}
\end{equation*}
这个公式表示度规微扰在沿矢量场 $\epsilon\xi^{\mu}$ 的无穷小微分同胚下的改变；我们将把它称为线性化理论中的\textbf{规范变换}（gauge transformation）。

微分同胚 $\psi_{\epsilon}$ 给出了同一物理情形的不同表示，同时保持我们对微扰很小的要求。因此，结果 (7.12) 告诉我们，什么样的度规微扰描写的时空在物理上等价——即对某个矢量 $\xi^{\mu}$ 而言彼此相差 $2\epsilon\partial_{(\mu}\xi_{\nu)}$ 的那些时空。我们的理论在这种变换下的不变性，类似于电磁学中在 $A_{\mu} \to A_{\mu} + \partial_{\mu}\lambda$ 之下的传统规范不变性。（这个类比不同于我们在附录 J 中与电磁学画的另一个类比；那个类比把正交标架形式中的局部洛伦兹变换联系到内部矢量丛中基的改变。）在电磁学中，不变性之所以成立，是因为场强 $F_{\mu\nu} = \partial_{\mu}A_{\nu} - \partial_{\nu}A_{\mu}$ 在规范变换下保持不变；类似地，我们发现变换 (7.14) 使线性化 Riemann 张量改变
\begin{align*}
\delta R_{\mu\nu\rho\sigma} = \frac{1}{2}(&\partial_{\rho}\partial_{\nu}\partial_{\mu}\xi_{\sigma} + \partial_{\rho}\partial_{\nu}\partial_{\sigma}\xi_{\mu} + \partial_{\sigma}\partial_{\mu}\partial_{\nu}\xi_{\rho} + \partial_{\sigma}\partial_{\mu}\partial_{\rho}\xi_{\nu} \\
&- \partial_{\sigma}\partial_{\nu}\partial_{\mu}\xi_{\rho} - \partial_{\sigma}\partial_{\nu}\partial_{\rho}\xi_{\mu} - \partial_{\rho}\partial_{\mu}\partial_{\nu}\xi_{\sigma} - \partial_{\rho}\partial_{\mu}\partial_{\sigma}\xi_{\nu}) \\
&= 0.
\tag{7.15}
\end{align*}
我们对度规微扰的适当规范变换所作的抽象推导，由此得到印证：它使曲率（从而也使物理时空）保持不变。

规范不变性也可以从一条观点稍微“低”一些（lowbrow）、却直接得多的途径来理解，即无穷小坐标变换。我们的微分同胚 $\psi_{\epsilon}$ 可以看作是把坐标从 $x^{\mu}$ 变到 $x^{\mu} - \epsilon\xi^{\mu}$。（这个不合常规的负号来自如下事实：“新”度规是从沿积分曲线向前一小段距离处拉回的，而这等价于把坐标换成沿曲线向后一小段距离处的坐标。）按照坐标变换下张量变换的通常规则一步步推下去，你可以恰好导出 (7.14)——尽管你得稍微作点弊，把两个不同坐标系中张量的分量等同起来。

% ---- 书页 279 / p294 ----
\section{自由度}

有了线性化 Einstein 张量的表达式 (7.8) 和规范变换效应的表达式 (7.14)，我们本可以立刻着手选取规范并求解 Einstein 方程。不过，我们可以先在 Minkowski 背景时空中选定一个固定的惯性坐标系，把度规微扰的分量按照它们在空间转动下的变换性质加以分解，由此积累起一些额外的物理洞见。你也许会担心，这样的分解有悖于广义相对论的坐标无关精神，但它其实与把电磁场强张量分解成电场和磁场并没有什么不同。尽管 $\mathbf{E}$ 和 $\mathbf{B}$ 都是一个 (0, 2) 张量的分量，但有时不妨假想某位固定观者的角色，把它们当作三维矢量来看待。\footnote{这里的讨论沿循 E. Bertschinger, “Cosmological Dynamics,” a talk given at Summer School on Cosmology and Large Scale Structure (Session 60), Les Houches, France, 1--28 Aug 1993; http://arXiv.org/abs/astro-ph/9503125。Bertschinger 关注的是宇宙学微扰论，其中类空超曲面随时间膨胀，但把它特化到不膨胀宇宙的情形也很简单。}

度规微扰是一个 (0, 2) 张量，但它不是反对称的，而是对称的。在空间转动下，00 分量是标量，0i 分量（等于 i0 分量）构成一个三维矢量，ij 分量则构成一个两指标的对称空间张量。这个空间张量还可以进一步分解成迹与无迹两部分。（用群论的语言说，我们在寻找转动群的“不可约表示”。换句话说，我们把张量分解成一个个单独的部分，它们在空间转动下只变换成它们自身。）于是我们把 $h_{\mu\nu}$ 写成
\begin{align*}
h_{00} &= -2\Phi \\
h_{0i} &= w_{i} \\
h_{ij} &= 2s_{ij} - 2\Psi\delta_{ij},
\tag{7.16}
\end{align*}
其中 $\Psi$ 编码了 $h_{ij}$ 的迹，而 $s_{ij}$ 是无迹的：
\begin{align*}
\Psi &= -\frac{1}{6}\delta^{ij}h_{ij} \\
s_{ij} &= \frac{1}{2}\left(h_{ij} - \frac{1}{3}\delta^{kl}h_{kl}\delta_{ij}\right).
\tag{7.17}
\end{align*}
整个度规于是可以写成
\begin{equation*}
\boxed{\,\mathrm{d}s^{2} = -(1+2\Phi)\mathrm{d}t^{2} + w_{i}(\mathrm{d}t\,\mathrm{d}x^{i} + \mathrm{d}x^{i}\,\mathrm{d}t) + [(1-2\Psi)\delta_{ij} + 2s_{ij}]\mathrm{d}x^{i}\mathrm{d}x^{j}.\,}
\tag{7.18}
\end{equation*}

% ---- 书页 280 / p295 ----
我们还没有选取规范，也没有求解任何方程，只是定义了一些方便的记号。无迹张量 $s_{ij}$ 被称为应变（strain），我们将会看到，引力辐射就包含在它之中。有时把空间分量分解成迹与无迹两部分并无帮助，这时我们可以径直坚持用 $h_{ij}$；在具体情形中哪种记号合适就用哪种。注意，正如在第 1 章中那样，空间度规现在就是简单的 $\delta_{ij}$，我们可以随意升降空间指标而不改变分量。

为了体会度规微扰中各个不同场的物理诠释，我们来考虑由测地线方程描述的检验粒子的运动。式 (7.18) 相应的 Christoffel 符号是
\begin{align*}
\Gamma^{0}_{00} &= \partial_{0}\Phi \\
\Gamma^{i}_{00} &= \partial_{i}\Phi + \partial_{0}w_{i} \\
\Gamma^{0}_{j0} &= \partial_{j}\Phi \\
\Gamma^{i}_{j0} &= \partial_{[j}w_{i]} + \frac{1}{2}\partial_{0}h_{ij} \\
\Gamma^{0}_{jk} &= -\partial_{(j}w_{k)} + \frac{1}{2}\partial_{0}h_{jk} \\
\Gamma^{i}_{jk} &= \partial_{(j}h_{k)i} - \frac{1}{2}\partial_{i}h_{jk}.
\tag{7.19}
\end{align*}
在这些表达式里，我们坚持用 $h_{ij}$ 而不用 $s_{ij}$ 和 $\Psi$，因为它们只是以组合 $h_{ij} = 2s_{ij} - 2\Psi\delta_{ij}$ 的形式进入。一旦我们开始取迹以得到 Ricci 张量和 Einstein 方程，这种区分就会变得重要。既然我们已经固定了一个惯性系，方便的做法是把四维动量 $p^{\mu} = \mathrm{d}x^{\mu}/\mathrm{d}\lambda$（若粒子有质量，则 $\lambda = \tau/m$）用能量 $E$ 和三速度 $v^{i} = \mathrm{d}x^{i}/\mathrm{d}t$ 表出：
\begin{equation*}
p^{0} = \frac{\mathrm{d}t}{\mathrm{d}\lambda} = E, \qquad p^{i} = Ev^{i}.
\tag{7.20}
\end{equation*}
然后我们取测地线方程
\begin{equation*}
\frac{\mathrm{d}p^{\mu}}{\mathrm{d}\lambda} + \Gamma^{\mu}_{\rho\sigma}p^{\rho}p^{\sigma} = 0,
\tag{7.21}
\end{equation*}
把第二项移到右边，使它呈现出力项的模样，再用 $E$ 除两边，就得到
\begin{equation*}
\frac{\mathrm{d}p^{\mu}}{\mathrm{d}t} = -\Gamma^{\mu}_{\rho\sigma}\frac{p^{\rho}p^{\sigma}}{E}.
\tag{7.22}
\end{equation*}
$\mu = 0$ 分量描述能量的演化，
\begin{equation*}
\frac{\mathrm{d}E}{\mathrm{d}t} = -E\left[\partial_{0}\Phi + 2(\partial_{k}\Phi)v^{k} - \left(\partial_{(j}w_{k)} - \frac{1}{2}\partial_{0}h_{jk}\right)v^{j}v^{k}\right].
\tag{7.23}
\end{equation*}

% ---- 书页 281 / p296 ----
你也许会认为能量应该守恒，但 $E = p^{0} = m\gamma$ 只包含粒子的“惯性”能量——在缓慢运动的极限下，即静能与动能——而不包含与引力场相互作用的那部分能量。

测地线方程的空间分量 $\mu = i$ 则变成
\begin{equation*}
\frac{\mathrm{d}p^{i}}{\mathrm{d}t} = -E\left[\partial_{i}\Phi + \partial_{0}w_{i} + 2(\partial_{[i}w_{j]} + \partial_{0}h_{ij})v^{j} + \left(\partial_{(j}h_{k)i} - \frac{1}{2}\partial_{i}h_{jk}\right)v^{j}v^{k}\right].
\tag{7.24}
\end{equation*}
为了作出物理解释，方便的做法是定义“引力电”（gravito-electric）与“引力磁”（gravito-magnetic）三维矢量场
\begin{align*}
G^{i} &\equiv -\partial_{i}\Phi - \partial_{0}w_{i} \\
H^{i} &\equiv (\nabla \times \vec{w})^{i} = \epsilon^{ijk}\partial_{j}w_{k},
\tag{7.25}
\end{align*}
它们与借助标势和矢势定义的普通电场和磁场有着显而易见的相似之处。于是 (7.24) 变成
\begin{equation*}
\frac{\mathrm{d}p^{i}}{\mathrm{d}t} = E\left[G^{i} + (\vec{v} \times H)^{i} - 2(\partial_{0}h_{ij})v^{j} - \left(\partial_{(j}h_{k)i} - \frac{1}{2}\partial_{i}h_{jk}\right)v^{j}v^{k}\right].
\tag{7.26}
\end{equation*}
右边的头两项描述的是，沿测地线运动的检验粒子如何响应标量微扰 $\Phi$ 和矢量微扰 $w_{i}$，其方式让人想起电磁学中的洛伦兹力定律。我们还发现与空间微扰 $h_{ij}$ 的耦合，它们对三速度分别是线性阶和二次阶的。各种微扰的相对重要性当然取决于所考虑的物理情形，我们很快就会展示这一点。

除了检验粒子的运动之外，我们还应当考察度规微扰的场方程，它们当然就是线性化 Einstein 方程。用我们的变量表出的 Riemann 张量是
\begin{align*}
R_{0j0l} &= \partial_{j}\partial_{l}\Phi + \partial_{0}\partial_{(j}w_{l)} - \frac{1}{2}\partial_{0}\partial_{0}h_{jl} \\
R_{0jkl} &= \partial_{j}\partial_{[k}w_{l]} - \partial_{0}\partial_{[k}h_{l]j} \\
R_{ijkl} &= \partial_{j}\partial_{[k}h_{l]i} - \partial_{i}\partial_{[k}h_{l]j},
\tag{7.27}
\end{align*}
其他分量由对称性相联系。我们用 $\eta^{\mu\nu}$ 缩并得到 Ricci 张量，
\begin{align*}
R_{00} &= \nabla^{2}\Phi + \partial_{0}\partial_{k}w^{k} + 3\partial_{0}^{2}\Psi \\
R_{0j} &= -\frac{1}{2}\nabla^{2}w_{j} + \frac{1}{2}\partial_{j}\partial_{k}w^{k} + 2\partial_{0}\partial_{j}\Psi + \partial_{0}\partial_{k}s_{j}{}^{k} \\
R_{ij} &= -\partial_{i}\partial_{j}(\Phi - \Psi) - \partial_{0}\partial_{(i}w_{j)} + \Box\Psi\delta_{ij} - \Box s_{ij} + 2\partial_{k}\partial_{(i}s_{j)}{}^{k},
\tag{7.28}
\end{align*}

% ---- 书页 282 / p297 ----
其中 $\nabla^{2} = \delta^{ij}\partial_{i}\partial_{j}$ 是三维平直空间中的拉普拉斯算符。由于 Ricci 张量涉及缩并，空间微扰的无迹部分与迹部分现在以不同的方式进入。最后，我们可以算出 Einstein 张量，
\begin{align*}
G_{00} &= 2\nabla^{2}\Psi + \partial_{k}\partial_{l}s^{kl} \\
G_{0j} &= -\frac{1}{2}\nabla^{2}w_{j} + \frac{1}{2}\partial_{j}\partial_{k}w^{k} + 2\partial_{0}\partial_{j}\Psi + \partial_{0}\partial_{k}s_{j}{}^{k} \\
G_{ij} &= (\delta_{ij}\nabla^{2} - \partial_{i}\partial_{j})(\Phi - \Psi) + \delta_{ij}\partial_{0}\partial_{k}w^{k} - \partial_{0}\partial_{(i}w_{j)} \\
&\qquad + 2\delta_{ij}\partial_{0}^{2}\Psi - \Box s_{ij} + 2\partial_{k}\partial_{(i}s_{j)}{}^{k} - \delta_{ij}\partial_{k}\partial_{l}s^{kl}.
\tag{7.29}
\end{align*}
把这个表达式用到 Einstein 方程 $G_{\mu\nu} = 8\pi GT_{\mu\nu}$ 中就会看出，度规分量中只有一小部分是引力场真正的自由度；其余的都遵守约束，由其他场来确定。为看清这一点，从 $G_{00} = 8\pi GT_{00}$ 出发，用 (7.29) 可以把它写成
\begin{equation*}
\nabla^{2}\Psi = 4\pi GT_{00} - \frac{1}{2}\partial_{k}\partial_{l}s^{kl}.
\tag{7.30}
\end{equation*}
这是一个关于 $\Psi$ 的、不含时间导数的方程；只要我们知道 $T_{00}$ 和 $s_{ij}$ 在任一时刻的行为，就能确定 $\Psi$ 必须是什么（相差空间无穷远处的边界条件）。因此，$\Psi$ 本身不是一个传播的自由度；它由能量-动量张量和引力应变 $s_{ij}$ 决定。接下来看 $0j$ 方程，我们把它写成
\begin{equation*}
(\delta_{jk}\nabla^{2} - \partial_{j}\partial_{k})w^{k} = -16\pi GT_{0j} + 4\partial_{0}\partial_{j}\Psi + 2\partial_{0}\partial_{k}s_{j}{}^{k}.
\tag{7.31}
\end{equation*}
这是一个关于 $w^{i}$ 的、不含时间导数的方程；同样地，只要我们知道能量-动量张量和应变（由它可以求出 $\Psi$），矢量 $w^{i}$ 就被确定了。最后，$ij$ 方程是
\begin{align*}
(\delta_{ij}\nabla^{2} - \partial_{i}\partial_{j})\Phi = 8\pi GT_{ij} + (\delta_{ij}\nabla^{2} - \partial_{i}\partial_{j} - 2\delta_{ij}\partial_{0}^{2})\Psi \\
-\ \delta_{ij}\partial_{0}\partial_{k}w^{k} + \partial_{0}\partial_{(i}w_{j)} + \Box s_{ij} - 2\partial_{k}\partial_{(i}s_{j)}{}^{k} - \delta_{ij}\partial_{k}\partial_{l}s^{kl}.
\tag{7.32}
\end{align*}
我们再一次看到，没有时间导数作用在 $\Phi$ 上，因此 $\Phi$ 被确定为其他场的函数。

这样一来，Einstein 方程中唯一传播的自由度就是应变张量 $s_{ij}$ 中的那些；我们将会看到，正是用它们来描述引力波。$h_{\mu\nu}$ 的其他分量都由 $s_{ij}$ 和物质场确定——它们不需要单独的初始数据。在其他替代理论中，比如 4.8 节讨论的那些含有额外场或作用量中高阶项的理论，度规的其他分量也可能成为动力学变量。正如我们将在 7.4 节末尾简要讨论的，传播的张量场经量子化后会给出具有不同自旋的粒子，这取决于场在空间转动下的行为。

% ---- 书页 283 / p298 ----
这样，标量 $\Phi$ 和 $\Psi$ 将是自旋-0 的，矢量 $w_{i}$ 将是自旋-1 的，而张量 $s_{ij}$ 是自旋-2 的。在通常的 GR 中，只有自旋-2 的部分才是真正的粒子激发。

在上一节中我们展示了规范变换 $h_{\mu\nu} \to h_{\mu\nu} + \partial_{\mu}\xi_{\nu} + \partial_{\nu}\xi_{\mu}$ 如何由一个矢量场 $\xi^{\mu}$ 生成。从现在起，我们把式 (7.14) 中的参数 $\epsilon$ 取为 1，而把矢量场 $\xi^{\mu}$ 本身看作小量。在这样的变换下，各个度规微扰场按如下方式改变：
\begin{align*}
\Phi &\to \Phi + \partial_{0}\xi^{0} \\
w_{i} &\to w_{i} + \partial_{0}\xi^{i} - \partial_{i}\xi^{0} \\
\Psi &\to \Psi - \frac{1}{3}\partial_{i}\xi^{i} \\
s_{ij} &\to s_{ij} + \partial_{(i}\xi_{j)} - \frac{1}{3}\partial_{k}\xi^{k}\delta_{ij},
\tag{7.33}
\end{align*}
这你可以很容易地验证。正如在电磁学和其他规范理论中那样，不同的规范适用于不同的情形；这里我们列出几种常用的选择。

首先考虑\textbf{横向规范}（transverse gauge，即宇宙学中有时使用的共形牛顿规范或 Poisson 规范的推广）。横向规范与电磁学中的 Coulomb 规范 $\partial_{i}A^{i} = 0$ 关系密切。我们先把应变固定成空间横向的，
\begin{equation*}
\partial_{i}s^{ij} = 0,
\tag{7.34}
\end{equation*}
即选取 $\xi^{j}$ 使之满足
\begin{equation*}
\nabla^{2}\xi^{j} + \frac{1}{3}\partial_{j}\partial_{i}\xi^{i} = -2\partial_{i}s^{ij}.
\tag{7.35}
\end{equation*}
$\xi^{0}$ 的值尚未确定，于是我们可以利用这个剩余的自由度，使矢量微扰成为横向的，
\begin{equation*}
\partial_{i}w^{i} = 0,
\tag{7.36}
\end{equation*}
即选取 $\xi^{0}$ 使之满足
\begin{equation*}
\nabla^{2}\xi^{0} = \partial_{i}w^{i} + \partial_{0}\partial_{i}\xi^{i}.
\tag{7.37}
\end{equation*}
作了傅里叶变换之后，“横向”的含义就清楚了：散度为零意味着张量与波矢量正交。(7.35) 和 (7.37) 都没有完全确定 $\xi^{\mu}$ 的值；它们都是含空间导数的二阶微分方程，需要边界条件才能确定解。就我们当前的目的而言，知道解总是存在的就足够了。条件 (7.34) 与 (7.36) 合起来定义了横向规范。在这个规范下，Einstein 方程变成
\begin{equation*}
\boxed{\,G_{00} = 2\nabla^{2}\Psi = 8\pi GT_{00},\,}
\tag{7.38}
\end{equation*}

% ---- 书页 284 / p299 ----
\begin{equation*}
\boxed{\,G_{0j} = -\frac{1}{2}\nabla^{2}w_{j} + 2\partial_{0}\partial_{j}\Psi = 8\pi GT_{0j},\,}
\tag{7.39}
\end{equation*}
以及
\begin{equation*}
\boxed{\,G_{ij} = (\delta_{ij}\nabla^{2} - \partial_{i}\partial_{j})(\Phi - \Psi) - \partial_{0}\partial_{(i}w_{j)} + 2\delta_{ij}\partial_{0}^{2}\Psi - \Box s_{ij} = 8\pi GT_{ij}.\,}
\tag{7.40}
\end{equation*}
在本章剩下的部分里，我们将利用这些方程在不同情形下求弱场解。

另一种常用的规范称为\textbf{同步规范}（synchronous gauge）。它等价于附录 D 中讨论的高斯法坐标的选取。可以把它看作电磁学中时间规范（temporal gauge）$A^{0} = 0$ 的引力类比，因为它把微扰的非空间分量统统消去。我们先令标势 $\Phi$ 消失，
\begin{equation*}
\Phi = 0,
\tag{7.41}
\end{equation*}
即选取 $\xi^{0}$ 使之满足
\begin{equation*}
\partial_{0}\xi^{0} = -\Phi.
\tag{7.42}
\end{equation*}
这之后我们仍然有能力选取 $\xi^{i}$。我们可以令矢量分量为零，
\begin{equation*}
w^{i} = 0,
\tag{7.43}
\end{equation*}
即选取 $\xi^{i}$ 使之满足
\begin{equation*}
\partial_{0}\xi^{i} = -w^{i} + \partial_{i}\xi^{0}.
\tag{7.44}
\end{equation*}
于是，同步规范下的度规呈现出如下诱人的形式：
\begin{equation*}
\mathrm{d}s^{2} = -\mathrm{d}t^{2} + (\delta_{ij} + h_{ij})\mathrm{d}x^{i}\mathrm{d}x^{j}.
\tag{7.45}
\end{equation*}
这只不过是规范选择的问题，而且适用于任何稍稍偏离 Minkowski 时空的时空。写出同步规范下的 Einstein 方程并不难，但我们不打算费这个事，因为本章剩下的部分实际上并不会用到它。

除了横向规范和同步规范之外，在计算引力波的产生时，方便的做法是用第三种选择，即 Lorenz 规范/谐和规范（Lorenz/harmonic gauge）。正如我们下面将要讨论的，它等价于令
\begin{equation*}
\partial_{\mu}h^{\mu}{}_{\nu} - \frac{1}{2}\partial_{\nu}h = 0,
\tag{7.46}
\end{equation*}

% ---- 书页 285 / p300 ----
其中 $h = \eta^{\mu\nu}h_{\mu\nu}$。用我们分解出的微扰场来表示，这个规范并没有任何特别简单的形式，但它确实使线性化 Einstein 方程呈现出一种格外简单的样子。

在继续转向弱场极限的应用之前，我们来结束关于自由度的讨论，着重指出 (7.16) 中对度规微扰分量的\emph{代数}分解，与另一种分解之间的区别——如果我们考虑的是张量\emph{场}而不是定义在一点上的张量，后一种分解就成为可能。这另一种分解有助于更直接地显现物理自由度，在宇宙学微扰论中至关重要。它的基础是一个标准的观察结果：一个矢量场可以分解成横向部分 $w_{\perp}^{i}$ 与纵向部分 $w_{\parallel}^{i}$：
\begin{equation*}
w^{i} = w_{\perp}^{i} + w_{\parallel}^{i},
\tag{7.47}
\end{equation*}
其中横向矢量是无散的，而纵向矢量是无旋的，
\begin{equation*}
\partial_{i}w_{\perp}^{i} = 0, \qquad \epsilon^{ijk}\partial_{j}w_{\parallel k} = 0.
\tag{7.48}
\end{equation*}
注意这些是微分方程，因此显然只有作用在张量场上才有意义。横向矢量可以表示成某个其他矢量 $\xi^{i}$ 的旋度，尽管 $\xi^{i}$ 的选取并不唯一，除非我们施加诸如 $\partial_{i}\xi^{i} = 0$ 这样的附加条件。纵向矢量则是标量 $\lambda$ 的散度，
\begin{equation*}
w_{\perp}^{i} = \epsilon^{ijk}\partial_{j}\xi_{k}, \qquad w_{\parallel i} = \partial_{i}\lambda.
\tag{7.49}
\end{equation*}
正如我们起初把度规微扰分解成标量、矢量和张量三部分一样，把矢量场分解成依赖于一个标量和一个横向矢量的两部分，这种分解在空间转动下是不变的。标量 $\lambda$ 显然代表一个自由度；矢量 $\xi^{i}$ 看起来像三个自由度，但由于 $\xi^{i}$ 选取的非唯一性，其中一个是虚的（你会注意到，这等价于作规范变换 $\xi_{i} \to \xi_{i} + \partial_{i}\omega$ 的自由）。因此总共有三个自由度，正好是描述原来的矢量场 $w^{i}$ 所应有的数目。

类似的步骤也适用于无迹对称张量 $s^{ij}$，它可以分解成横向部分 $s_{\perp}^{ij}$、螺线部分 $s_{\mathrm{S}}^{ij}$ 和纵向部分 $s_{\parallel}^{ij}$，
\begin{equation*}
s^{ij} = s_{\perp}^{ij} + s_{\mathrm{S}}^{ij} + s_{\parallel}^{ij}.
\tag{7.50}
\end{equation*}
横向部分是无散的，螺线部分的散度是一个横向（无散）矢量，而纵向部分的散度是一个纵向（无旋）矢量：
\begin{align*}
\partial_{i}s_{\perp}^{ij} = 0 \\
\partial_{i}\partial_{j}s_{\mathrm{S}}^{ij} = 0 \\
\epsilon^{jkl}\partial_{k}\partial_{i}s_{\parallel}{}^{i}{}_{j} = 0.
\tag{7.51}
\end{align*}

% ---- 书页 286 / p301 ----
这意味着纵向部分可以由一个标量场 $\theta$ 导出，而螺线部分可以由一个横向矢量 $\zeta^{i}$ 导出，
\begin{align*}
s_{\parallel ij} &= \left(\partial_{i}\partial_{j} - \frac{1}{3}\delta_{ij}\nabla^{2}\right)\theta \\
s_{\mathrm{S}ij} &= \partial_{(i}\zeta_{j)},
\tag{7.52}
\end{align*}
其中
\begin{equation*}
\partial_{i}\zeta^{i} = 0.
\tag{7.53}
\end{equation*}
于是，纵向部分描写一个单独的自由度，而螺线部分描写两个自由度。横向部分无法再进一步分解；它描写对称无迹的 $3 \times 3$ 张量 $s_{ij}$ 其余的两个自由度。在本章后面，我们将引入横无迹（TT）规范来描述真空中传播的引力波；在这个规范下，唯一非零的度规微扰就是横向张量微扰 $s_{\perp}^{ij}$。

借助这种对张量场的分解，我们成功地把原先有十个分量的度规微扰 $h_{\mu\nu}$ 写成了四个标量（$\Phi$、$\Psi$、$\lambda$ 和 $\theta$），各带一个自由度；两个横向矢量（$\xi^{i}$ 和 $\zeta^{i}$），各带两个自由度；以及一个横无迹张量（$s_{\perp}^{ij}$），带两个自由度。当人们谈论“标量”“矢量”和“张量”模式时，指的就是这组场。接下来我们可以用类似的方式分解能量-动量张量，用这些变量写出 Einstein 方程，并分离出物理的（规范不变的）自由度。本书不会使用这种分解，但你在查阅文献时应当知道它的存在。

\section{Newton 场与光子轨迹}

我们此前把“Newton 极限”定义为描述源静态、检验粒子缓慢运动的弱场。本节我们将稍稍扩展这个定义：仍然只限于静态的源，但允许检验粒子以任意速度运动。这里显然有一个重要差别：我们此前只需要考虑度规 $g_{00}$ 分量的效应，但我们将会发现，相对论性粒子对度规的空间分量也有响应。

我们可以用尘埃——即压强为零的理想流体——来为静态的引力源建模。（宇宙中的大多数物质都可以很好地用尘埃来近似，包括恒星、行星、星系，甚至暗物质。）我们在
% 本块末句在原书书页286末尾截断，于原书书页287顶部继续（“frame of the dust, ...”）；下一块 ch7_b 已以 %JOIN% 续接本句。
